\scanpage{41}
opère trivialt, \inserted{(et réciproquement bien sûr)} Donc \unsure{du pt de} vue virtuel, \inserted{et pour $U$ fixé,} les modules $E$ envisagés sont \unsure{essentiellement} les \unsure{$=$s|mêmes} que les modules tels qu'\deleted{il existe un ss-grp \unsure{tel que}} $U$ opère trivialt sur $E$.

Or pour un $E$ sur lequel $U$ opère trivialt, on peut prendre l'opération moyenne
\[
E\leadsto E^I\qquad\text{de}\quad \operatorname{Mod}(D/U,k)-\operatorname{Mod}(D/I,k)
\]
et si $k$ est de car. $0$, c'est là un \deleted{opération \uline{additive}, qui \unsure{nous donne}} \inserted{foncteur exact. Il induit un hom.} \deleted{\unsure{de gr.}} virtuel \deleted{$\unsure{E}$} $E\leadsto E^{\unsure{\natural|\xi}}$
\[
\unsure{\natural|\xi}_U:\ R(D/U,k)\longrightarrow R(D/I,k)
\]
On a. compatibilité, pour $V\subset U$, pour les \uline{foncteurs} $E\leadsto E^I$, d'où pour les $\unsure{\natural|\xi}_U$ et $\unsure{\natural|\xi}_V$ :
\diagram{\begin{tikzcd}[column sep=large,row sep=large]
R(D/U,k) \arrow[r,no head,"{\unsure{\natural|\xi}_U}"] \arrow[d,"{\text{\uline{\unsure{restr.}}}\ \deleted{\unsure{\natural}}}"'] & R(D/I,k) \\
R(D/V,k) \arrow[ur,"{\unsure{\natural|\xi}_V}"'] &
\end{tikzcd}}

\margin{40}

